\phantom{\cite{fukushima1980neocognitron, wiskott2002slow, hinton2006fast,oquab2014learning, simonyan2014very, girshick2014rich, long2015fully}}

\begin{wrapfigure}{r}{0.50\textwidth}
  \vspace{-7.7mm}
   \centering
   \begin{minipage}{0.5\textwidth}
    \centering
     \includegraphics[width=0.95\textwidth,bb=0 0 1140 1018]{img/byol.pdf}
       % \includegraphics[width=0.95\textwidth]{img/byol.pdf}
        \caption{Performance of \algo on ImageNet (linear evaluation) using ResNet-50 and our best architecture ResNet-200 ($2\times$), compared to other unsupervised and supervised (\texttt{Sup\!.\!}) baselines~\cite{SimCLR}.}
        \label{fig:holygrail}
 \end{minipage}
 \vspace{-8mm}
\end{wrapfigure}

Learning good image representations is a key challenge in computer vision~\cite{fukushima1980neocognitron, wiskott2002slow, hinton2006fast} as it allows for efficient training on downstream tasks~\cite{oquab2014learning, simonyan2014very, girshick2014rich, long2015fully}.
%
Many different training approaches have been proposed to learn such representations, usually relying on visual pretext tasks. 
Among them, state-of-the-art contrastive methods~\cite{SimCLR, MOCO, CPC, CMC, Tian2020WhatMF} are trained by reducing the distance between representations of different \augview s of the same image (`positive pairs'), and increasing the distance between representations of \augview s from different images (`negative pairs'). These methods need careful treatment of negative pairs~\cite{pmlr-v97-saunshi19a} by either relying on large batch sizes~\cite{SimCLR, Tian2020WhatMF}, memory banks~\cite{MOCO} or customized mining strategies~\cite{Manmatha2017SamplingMI, Harwood2017SmartMF} to retrieve the negative pairs. In addition, their performance critically depends on the choice of \imageaugmentation s~\cite{SimCLR,Tian2020WhatMF}. 

In this paper, we introduce \longalgo~(\algo), a new algorithm for self-supervised learning of image representations. \algo achieves higher performance than state-of-the-art contrastive methods without using negative pairs. It iteratively bootstraps\footnote{Throughout this paper, the term \emph{bootstrap} is used in its idiomatic sense rather than the statistical sense.} the outputs of a network to serve as targets for an enhanced representation.
Moreover, \algo is more robust to the choice of \imageaugmentation s than contrastive methods; we suspect that not relying on negative pairs is one of the leading reasons for its improved robustness. While previous methods based on bootstrapping have used pseudo-labels~\cite{lee2013PseudoLabelT}, cluster indices~\cite{caron2018DeepCF} or a handful of labels~\cite{bachman2014learning, laine2016temporal, tarvainen2017mean}, 
we propose to \textit{directly bootstrap the representations}. In particular, \algo uses two neural networks, referred to as online and target networks, that interact and learn from each other. Starting from \aaugview~of an image, \algo trains its online network to predict the target network's representation of another \augview~of the same image. 
While this objective admits collapsed solutions, e.g., outputting the same vector for all images, we empirically show that \algo does not converge to such solutions. 
We hypothesize (see~\Cref{sec:intuition_byol}) that the combination of
(i) the addition of a predictor to the online network 
and (ii) the use of a slow-moving average of the online parameters as the target network encourages encoding more and more information within the online projection and avoids collapsed solutions.

We evaluate the representation learned by \algo on ImageNet~\cite{ILSVRC15} and other vision benchmarks using ResNet architectures~\cite{ResNet}. Under the linear evaluation protocol on ImageNet, consisting in training a linear classifier on top of the frozen representation, \algo reaches $74.3\%$ top-1 accuracy with a standard ResNet-$50$ and $79.6\%$ top-1 accuracy with a larger ResNet (\Cref{fig:holygrail}). 
In the semi-supervised and transfer settings on ImageNet, we obtain results on par or superior to the current state of the art. Our contributions are: (\textit{i}) We introduce \algo, a self-supervised representation learning method (\Cref{sec:method}) which achieves state-of-the-art results under the linear evaluation protocol on ImageNet without using negative pairs. (\textit{ii}) We show that our learned representation outperforms the state of the art on semi-supervised and transfer benchmarks (\Cref{sec:experiments}). (\textit{iii}) We show that \algo is more resilient to changes in the batch size and in the set of \imageaugmentation s compared to its contrastive counterparts (\Cref{sec:ablation}).
In particular, \algo suffers a much smaller performance drop than \simclr, a strong contrastive baseline, when only using random crops as \imageaugmentation s.
